% 既知の制約と拡張時の課題
\section{既知の制約と拡張時の課題}

本アーキテクチャは小〜中規模（モジュール10個程度まで）のプロジェクトに適した設計であるが、
システムの複雑化・大規模化に伴い、以下の制約が顕在化する可能性がある。

\subsection{SystemDataの肥大化}

モジュールが増えるたびに\texttt{SystemData}にサブ構造体が追加され、構造体が肥大化する。

\begin{itemize}
    \item 全モジュールが\texttt{SystemData\&}を受け取るため、自身に無関係なデータにもアクセス可能になる。
    入力モジュールが出力用データを書き換える等の誤用をコンパイラレベルで防止できない
    \item モジュール1つの追加でも\texttt{ProjectConfig.h}の変更が必要となり、全モジュールの再コンパイルが発生する
    \item モジュール数が15〜20を超えると、構造体の見通しが悪化し保守性が低下する
\end{itemize}

\subsection{モジュール間の依存関係}

現在のアーキテクチャでは、モジュール間の初期化順序や実行順序の依存関係を明示的に表現する仕組みがない。

\begin{itemize}
    \item \texttt{init()}の呼び出し順序は配列の定義順に暗黙的に依存しており、コード上の意図が見えにくい
    \item 例：カメラモジュールがSDカードモジュールの初期化完了を前提とする場合、配列の順番に依存する脆い設計になる
    \item モジュール数が増えると、暗黙の順序依存が複雑化し、並べ替えによるバグが発生しやすくなる
\end{itemize}

\subsection{リアルタイム性の限界}

3フェーズモデルは処理の流れが明快であるが、1ループ周期分の応答遅延が必ず発生する。

\begin{itemize}
    \item 入力フェーズで検知した異常に対して、同一フレーム内で出力モジュールが即座に反応できない
    （ロジックフェーズ→出力フェーズを経由する必要がある）
    \item 緊急停止のようなリアルタイム性が要求されるケースでは、1フレームの遅延が致命的になりうる
    \item 割り込み（ISR）との統合パターンが定義されておらず、ハードウェア割り込みを活用した即時応答が困難
\end{itemize}

\subsection{シングルスレッド制約}

本アーキテクチャは協調的なシングルスレッドスケジューリングに基づいており、
ESP32-S3のデュアルコアを活用する仕組みを持たない。

\begin{itemize}
    \item カメラ画像処理やFFT等の重い処理が1つのモジュールに含まれると、全モジュールの更新が遅延する
    \item FreeRTOSタスクとの共存パターンが未定義のため、WiFi通信やBLE処理等の非同期I/Oとの統合が難しい
    \item \texttt{update()}の処理時間に上限を設ける仕組みがなく、モジュール数の増加に伴い1ループの累積実行時間が問題になる
\end{itemize}

\subsection{ロジックフェーズのスケーラビリティ}

\texttt{applyPattern()}等の関数でロジックを記述する方式は、モジュール数の増加に伴い複雑化する。

\begin{itemize}
    \item 入力モジュールと出力モジュールの組み合わせが増えると、ロジック関数が肥大化する
    \item 複数のロジックが同一の出力データを書き換える場合の優先度・排他制御の仕組みがない
    \item 例：温度センサーAと温度センサーBの両方がモーターの速度を制御したい場合、競合の解決方法が未定義
\end{itemize}

\subsection{動的なモジュール管理}

モジュール配列は\texttt{setup()}時に静的に確定され、実行中の変更ができない。

\begin{itemize}
    \item 実行中にモジュールを有効化・無効化する仕組みがない
    （例：BLE経由の設定変更でモジュールを切り替える）
    \item ホットプラグ的なハードウェア（USB接続デバイス等）への対応が困難
\end{itemize}

\subsection{エラー復帰パターンの不在}

エラーの検出手段は提供されているが、エラーからの復帰に関するパターンが未定義である。

\begin{itemize}
    \item \texttt{init()}が\texttt{false}を返した場合のリトライ、スキップ、システム停止等の判断基準が定義されていない
    \item \texttt{update()}でエラーフラグを設定した後、そのフラグを誰がいつクリアするかが不明確
    \item ウォッチドッグタイマーやフェイルセーフモードへの移行パターンが未定義
\end{itemize}

\subsection{Configの実行時変更}

Config構造体はコンストラクタで\texttt{const\&}としてコピーされ、以降の変更ができない。

\begin{itemize}
    \item BLEやMQTT経由でキャリブレーション値やしきい値を実行時に変更するユースケースに対応できない
    \item 変更にはモジュールの再構築が必要だが、実行中にモジュールを再構築する仕組みがない
\end{itemize}

\subsection{テスタビリティ}

ハードウェア依存が強く、PC上でのユニットテストが困難である。

\begin{itemize}
    \item \texttt{IModule<SystemData>}と具象型に束縛されるため、モック化が難しい
    \item \texttt{Arduino.h}や\texttt{millis()}等のハードウェア依存APIを抽象化する仕組みがない
    \item CI環境でのビルド検証は可能だが、ロジックの自動テストを行う基盤が整っていない
\end{itemize}

\begin{notebox}[本セクションの位置づけ]
ここに挙げた課題は、アーキテクチャの設計上の欠陥ではなく、
\textbf{簡潔さとのトレードオフ}として意図的に受け入れている制約である。\\
プロジェクトの規模や要件に応じて、イベントバス、タスク分離、依存グラフ等のパターンを
追加導入することで対処可能であるが、アーキテクチャの複雑化を伴うため慎重に判断すべきである。
\end{notebox}
